\documentclass[a4paper,11pt]{article}

\usepackage[utf8]{inputenc}
\usepackage[T1]{fontenc}
\usepackage[french]{babel}
\usepackage{geometry}
\geometry{margin=2.5cm}
\usepackage{hyperref}
\usepackage{listings}
\usepackage{xcolor}
\usepackage{booktabs}
\usepackage{array}

\lstdefinestyle{ccode}{language=C,basicstyle=\ttfamily\small,keywordstyle=\color{blue},
  commentstyle=\color{gray},stringstyle=\color{orange},breaklines=true,frame=single,
  numbers=left,numberstyle=\tiny\color{gray},tabsize=4,columns=fullflexible}
\lstdefinestyle{ktcode}{language=Java,basicstyle=\ttfamily\small,keywordstyle=\color{blue!80!black},
  commentstyle=\color{gray},stringstyle=\color{red!70!black},breaklines=true,frame=single,
  numbers=left,numberstyle=\tiny\color{gray},tabsize=4,columns=fullflexible}
\lstdefinestyle{bashcode}{language=bash,basicstyle=\ttfamily\small,keywordstyle=\color{blue},
  commentstyle=\color{gray},stringstyle=\color{orange},breaklines=true,frame=single,columns=fullflexible}
\lstdefinestyle{confcode}{basicstyle=\ttfamily\small,breaklines=true,frame=single,
  backgroundcolor=\color{gray!10},columns=fullflexible}

\title{\textbf{Documentation technique --- Mini Dashboard}\\
\large Tableau de bord LVGL sur Raspberry Pi~5 (VERSION COMPLeTE)\\
\normalsize \texttt{lv\_port\_vscode/src} \quad | \quad \texttt{NavRelay} \quad | \quad \texttt{SETUP.sh}}
\author{meka3}
\date{\today}

\begin{document}
\maketitle
\tableofcontents
\newpage

%% ===========================================================================
\section{DÉMARRAGE RAPIDE (NOUVELLE INSTALLATION)}
%% ===========================================================================

\subsection{Installation en 3 etapes}

Pour une nouvelle installation de Raspberry Pi OS Bookworm, executer :

\begin{enumerate}
  \item \textbf{Cloner le depôt et les sous-modules}
  \begin{lstlisting}[style=bashcode]
git clone --recursive --branch master https://github.com/EEmeka33/Mini_Dashboard.git
cd Mini_Dashboard
  \end{lstlisting}
  \emph{Note:} L'option \texttt{--recursive} remplace les deux commandes separees
  \texttt{git submodule update --init --recursive}, simplifiant le processus.
  
  \item \textbf{Installation et configuration automatisée (SETUP.sh)}
  \begin{lstlisting}[style=bashcode]
bash SETUP.sh
  \end{lstlisting}
  Duree: ~5-10 minutes. Installe paquets, configure services systemd et le daemon RFCOMM.
  
  Le script installe et active automatiquement :
\begin{itemize}
    \item \texttt{bluetooth.service} pour BlueZ
    \item \texttt{ofono.service} pour les appels HFP
    \item \texttt{nav-bind.service} pour le daemon RFCOMM (nouvelle feature)
    \item \texttt{pulseaudio} pour la musique Bluetooth
\end{itemize}
  
  \item \textbf{Compiler le dashboard}
  \begin{lstlisting}[style=bashcode]
cd lv_port_pc_vscode
rm -rf build
mkdir build
cd build
cmake ..
make -j4
  \end{lstlisting}
  
  \item \textbf{Vérifier que les services système tournent}
  \begin{lstlisting}[style=bashcode]
sudo systemctl status bluetooth.service
sudo systemctl status ofono.service
sudo systemctl status nav-bind.service
  \end{lstlisting}
  Tous les trois doivent montrer \texttt{active (running)}.

  \item \textbf{Lancer le fichier compilé :}
  \begin{lstlisting}[style=bashcode]
cd lv_port_pc_vscode/bin/
./main
  \end{lstlisting}
  
\end{enumerate}


\newpage

%% ===========================================================================
\section{Vue d'ensemble du système}
%% ===========================================================================

Le projet Mini Dashboard est un tableau de bord automobile embarqué tournant sur Raspberry Pi~5
Sous un écran circulaire $1080\times1080$ pixels. Il agrège trois flux de données en temps réel :

\begin{itemize}
  \item \textbf{Navigation} : données de Google Maps relayées depuis un smartphone Android via
    Bluetooth RFCOMM (profil SPP), lues sur \texttt{/dev/rfcomm1} sous forme de lignes JSON.
  \item \textbf{Musique Bluetooth} : métadonnées A2DP/AVRCP et jaquette, obtenues par D-Bus via
    \texttt{org.bluez.MediaPlayer1}.
  \item \textbf{Appels HFP} : détection et gestion des appels entrants via oFono et
    \texttt{org.ofono.VoiceCallManager}.
\end{itemize}

\subsection{Architecture globale}

\begin{center}
\begin{tabular}{>{\centering\arraybackslash}m{3.8cm}|>{\centering\arraybackslash}m{4.2cm}|>{\centering\arraybackslash}m{4.2cm}}
\toprule
\textbf{Smartphone Android} & \textbf{Couche systeme Pi} & \textbf{Application LVGL} \\
\midrule
NavRelay (Kotlin) & BlueZ + RFCOMM & \texttt{nav\_thread\_fn()} \\
Google Maps notifications & \texttt{/dev/rfcomm1} & JSON $\to$ \texttt{nav\_info\_t} \\
\hline
Lecteur audio A2DP/AVRCP & BlueZ \texttt{MediaPlayer1} & \texttt{music\_thread\_fn()} \\
& D-Bus systeme & \texttt{get\_music\_properties()} \\
\hline
Telephonie Android & BlueZ + oFono & \texttt{calls\_thread\_fn()} \\
Profil HFP & D-Bus systeme & \texttt{org.ofono.VoiceCall} \\
\bottomrule
\end{tabular}
\end{center}

\subsection{Prérequis matériels et logiciels}

\begin{itemize}
  \item Raspberry Pi~5 avec Raspberry Pi OS Bookworm (64-bit recommande)
  \item ecran $1080\times1080$, connecte via HDMI ou DSI
  \item Contrôleur Bluetooth integre ou USB (\texttt{hci0})
  \item Smartphone Android avec l'application \textbf{NavRelay} installee et Google Maps
\end{itemize}

\subsubsection{Build tools (obligatoire pour compiler)}

\begin{lstlisting}[style=bashcode]
sudo apt update && sudo apt install -y \
  build-essential cmake git pkg-config
\end{lstlisting}

\subsubsection{Paquets intime et dépendances}

\begin{lstlisting}[style=bashcode]
sudo apt update
sudo apt install -y pulseaudio pulseaudio-module-bluetooth
sudo apt install -y bluez bluez-tools ofono modem-manager
sudo apt install -y libdbus-1-dev libcurl4-openssl-dev libjson-c-dev
sudo apt install -y libsdl2-dev libsdl2-image-dev
sudo apt install -y alsa-utils dbus
\end{lstlisting}

Ou en une seule commande (identique à SETUP.sh) :

\begin{lstlisting}[style=bashcode]
sudo apt update
sudo apt install -y pulseaudio pulseaudio-module-bluetooth bluez
sudo apt install -y bluez-tools ofono modem-manager alsa-utils dbus
sudo apt install -y libdbus-1-dev libcurl4-openssl-dev libjson-c-dev
sudo apt install -y libsdl2-dev libsdl2-image-dev
\end{lstlisting}

\subsubsection{Installation automatisée (RECOMMANDÉE)}

Tous les paquets et configurations ci-dessus sont appliqués automatiquement :

\begin{lstlisting}[style=bashcode]
bash /home/pi5/dashboard/SETUP.sh
\end{lstlisting}

Ce script prend 5-10 minutes et configure également les services systemd dans les bonnes conditions.

%% ===========================================================================
\section{Compilation et mise en place du projet LVGL}
%% ===========================================================================

\subsection{Cloner le depôt et initialiser les sous-modules}

Le projet utilise LVGL comme sous-module Git. L'option \texttt{--recursive} initialise
automatiquement les sous-modules lors du clonage :

\begin{lstlisting}[style=bashcode]
git clone --recursive https://github.com/EEmeka33/Mini_Dashboard.git
cd Mini_Dashboard
\end{lstlisting}

\emph{Avantage :} Une seule commande \texttt{git clone --recursive} remplace les deux
commandes separees \texttt{git clone} puis \texttt{git submodule update --init --recursive}.

\newpage

\subsection{Structure du repertoire \texttt{lv\_port\_vscode}}

\begin{verbatim}
lv_port_vscode/
  src/
    hal/
        hal.c / hal.h       # Pilote SDL2 pour LVGL
    main.c                  # Point d'entree, UI, threads
    speedo_interactive.c    # Widget compteur de vitesse
    speedo_interactive.h    # API publique du compteur
    album_art.c             # Recuperation jaquette iTunes
    bt_pairing.h            # Appels Pair button
  lvgl/                     # Sources LVGL (sous-module)
  lv_conf.h                 # Configuration LVGL
  CMakeLists.txt
  build/                    # Repertoire de compilation
\end{verbatim}

\subsection{Configuration LVGL (\texttt{lv\_conf.h})}

Les polices suivantes doivent être activees car elles sont utilisees dans
\texttt{main.c} et \texttt{speedo\_interactive.c} :

\begin{lstlisting}[style=confcode,caption={lv\_conf.h --- polices a activer}]
#define LV_FONT_MONTSERRAT_14  1
#define LV_FONT_MONTSERRAT_16  1
#define LV_FONT_MONTSERRAT_18  1
#define LV_FONT_MONTSERRAT_22  1
#define LV_FONT_MONTSERRAT_24  1
#define LV_FONT_MONTSERRAT_48  1
\end{lstlisting}

Le pilote systeme de fichiers POSIX doit être active pour le chargement d'images JPEG :

\begin{lstlisting}[style=confcode]
#define LV_USE_FS_POSIX      1
#define LV_FS_POSIX_LETTER   'A'
\end{lstlisting}

Le decodeur JPEG integre doit être active pour que \texttt{lv\_image\_set\_src()} charge
\texttt{cover.jpg} :

\begin{lstlisting}[style=confcode]
#define LV_USE_TJPGD  1   /* ou LV_USE_LIBJPEG_TURBO selon la version */
\end{lstlisting}

Ligne pour afficher l'application en Fullscreen :

\begin{lstlisting}[style=confcode]
#define LV_SDL_FULLSCREEN       1
\end{lstlisting}


\subsection{Compilation complete avec CMake (RECOMMANDÉE)}

\begin{lstlisting}[style=bashcode]
cd lv_port_vscode
rm -rf build
mkdir build
cd build
cmake ..
make -j4
\end{lstlisting}

Ou en une seule ligne :

\begin{lstlisting}[style=bashcode]
cd lv_port_vscode && rm -rf build && mkdir build && cd build && cmake .. && make -j4
\end{lstlisting}

L'exécutable sera créé à \texttt{bin/main}.

\subsection{Compilation manuelle (alternative)}

\begin{lstlisting}[style=bashcode]
cd lv_port_vscode
gcc -O2 -o dashboard \
  src/main.c src/speedo_interactive.c src/album_art.c \
  $(find lvgl/src -name "*.c" | tr '\n' ' ') \
  hal/hal.c \
  $(pkg-config --cflags --libs dbus-1 libcurl json-c sdl2) \
  -I. -Ilvgl -Ilvgl/src \
  -lm -lpthread
\end{lstlisting}

\subsection{Lancer le dashboard}

Le programme necessite une session graphique (SDL2) :

\begin{lstlisting}[style=bashcode]
export DISPLAY=:0
cd lv_port_vscode/bin
./main
\end{lstlisting}

Pour un lancement automatique au demarrage, un service systemd peut être configure
via SETUP.sh.

%% ===========================================================================
\section{Pair Button et decouverte Bluetooth}
%% ===========================================================================

\subsection{Nouvelle approche simplifiee}

Au lieu de tenter un appairage automatise par script (qui conflitue avec l'interface
Bluetooth du Pi OS), le \textbf{Pair Button} utilise une approche deux-etapes :

\begin{enumerate}
  \item \textbf{Utilisateur clique sur le boutton "Pair"} dans le dashboard (UI LVGL).
  \item \textbf{Pi devient decouvrable} pendant 3 minutes via le script \texttt{bt\_make\_discoverable.sh}.
  \item \textbf{Utilisateur appaire depuis Bluetooth Settings du Pi} ou du telephone.
\end{enumerate}

Cela elimine entierement le conflit logiciel car une seule application (l'interface
native du Pi) envoie les commandes Bluetooth.

\subsection{Rôle de \texttt{bt\_pairing.h}}

\begin{lstlisting}[style=ccode,caption={bt\_pairing.h --- fonction publique}]
static inline int bluetooth_pair_device_start(void)
{
    /* Lance bt_make_discoverable.sh en processus fils detache */
    pid_t pid = fork();
    if (pid == 0) {
        setsid();  /* Detache de la session parente */
        execv("/bin/bash", ...);    
        exit(0);
    }
    return (pid > 0) ? 0 : -1;
}
\end{lstlisting}

La fonction retourne 0 en cas de succes, -1 en cas d'erreur. Le callback bouton
affiche \texttt{"Made Discoverable"} ou \texttt{"Error"}.

\subsection{Script \texttt{bt\_make\_discoverable.sh}}

\begin{lstlisting}[style=bashcode,caption={bt\_make\_discoverable.sh}]
#!/bin/bash
bluetoothctl discoverable on
sleep 180         # 3 minutes
bluetoothctl discoverable off
\end{lstlisting}

Ce script s'execute en arriere-plan, sans bloquer l'UI du dashboard.

\subsection{Appairage initial du telephone}

\subsubsection{Methode 1 : Via le boutton Pair du dashboard (RECOMMANDÉE)}

\begin{enumerate}
  \item \textbf{Lancer le dashboard} : \texttt{./bin/main} depuis \texttt{lv\_port\_pc\_vscode/build}
  \item \textbf{Cliquer sur le boutton bleu \og Pair\fg{}} dans l'UI LVGL
  \item \textbf{Le Pi devient decouvrable pendant 3 minutes}
  \item \textbf{Sur le telephone} : ouvrir Parametres $\to$ Bluetooth $\to$ Chercher des appareils
  \item \textbf{Selectionner \og raspberrypi\fg{}} et appuyer sur Appairer
  \item \textbf{Confirmer le PIN} sur le telephone si demande
  \item \textbf{Dashboard affiche \og Made Discoverable\fg{}} si succes
\end{enumerate}

\emph{Avantage:} Aucune ligne de commande, interface graphique integree au dashboard.
Le MAC du telephone est automatiquement detecte par \texttt{main.c}.

\subsubsection{Methode 2 : En ligne de commande (manuel)}

\begin{lstlisting}[style=bashcode]
sudo bluetoothctl
[bluetooth]# power on
[bluetooth]# agent on
[bluetooth]# default-agent
[bluetooth]# scan on
\end{lstlisting}

Attendre que le telephone apparaisse :

\begin{verbatim}
[NEW] Device 2C:BE:EB:7F:2D:FA Pixel 7
\end{verbatim}

Continuer l'appairage :

\begin{lstlisting}[style=bashcode]
[bluetooth]# scan off
[bluetooth]# pair 2C:BE:EB:7F:2D:FA
\end{lstlisting}

Confirmer le PIN sur le telephone (ou pas si le Pi envoie le PIN automatiquement).

\begin{lstlisting}[style=bashcode]
[bluetooth]# trust 2C:BE:EB:7F:2D:FA
[bluetooth]# connect 2C:BE:EB:7F:2D:FA
[bluetooth]# quit
\end{lstlisting}

\emph{Note:} Apres appairage, l'adresse MAC est automatiquement sauvegardee par
\texttt{main.c} dans \texttt{/tmp/connected\_phone\_mac.txt} lors de la detection
du telephone. Aucune modification manuelle du daemon n'est necessaire.

%% ===========================================================================
\section{Configuration systeme Bluetooth}
%% ===========================================================================

\subsection{Configuration BlueZ pour la musique et les appels}

\subsubsection{Desactivation du profil Headset natif de BlueZ}

Par defaut, BlueZ enregistre les UUID HFP/HSP pour son propre profil \emph{Headset},
ce qui entre en conflit avec oFono. L'erreur suivante apparaît alors :

\begin{verbatim}
RegisterProfile() replied: org.bluez.Error.NotPermitted
UUID already registered
\end{verbatim}

Pour ceder la gestion HFP à oFono, editer \texttt{/etc/bluetooth/main.conf} :

\begin{lstlisting}[style=bashcode]
sudo nano /etc/bluetooth/main.conf
\end{lstlisting}

\begin{lstlisting}[style=confcode,caption={/etc/bluetooth/main.conf}]
[General]
Disable=Headset
\end{lstlisting}

Cette directive ne desactive ni A2DP ni AVRCP ; seul le profil casque/mains libres
natif BlueZ est desactive. Le lecteur de musique continue de fonctionner normalement.

\subsubsection{Activation du service oFono}

\begin{lstlisting}[style=bashcode]
sudo systemctl enable ofono
sudo systemctl start ofono
sudo systemctl status ofono
\end{lstlisting}

\subsubsection{Redemarrage des services Bluetooth}

Apres toute modification de \texttt{main.conf} :

\begin{lstlisting}[style=bashcode]
sudo systemctl restart bluetooth
sudo systemctl restart ofono
\end{lstlisting}

\subsection{Configuration PipeWire pour HFP via oFono}

Sur Raspberry Pi OS Bookworm, PipeWire remplace PulseAudio. Pour que les appels HFP
passent par oFono :

\begin{lstlisting}[style=bashcode]
sudo mkdir -p /etc/pipewire/pipewire-pulse.conf.d/
sudo nano /etc/pipewire/pipewire-pulse.conf.d/10-ofono.conf
\end{lstlisting}

\begin{lstlisting}[style=confcode,caption={10-ofono.conf}]
pulse.properties = {
  bluez5.headset-roles = [ hfp_hf hsp_hs ]
  bluez5.hfphsp-backend = ofono
}
\end{lstlisting}

Redemarrer PipeWire dans la session utilisateur (sans \texttt{sudo}) :

\begin{lstlisting}[style=bashcode]
systemctl --user restart pipewire pipewire-pulse wireplumber
\end{lstlisting}

\subsection{Verification de la pile D-Bus}

Apres connexion du telephone et lancement d'un morceau :

\begin{lstlisting}[style=bashcode]
# Verifier la presence d'un objet MediaPlayer1 (musique)
busctl tree org.bluez
# Attendu : /org/bluez/hci0/dev_2C_BE_EB_7F_2D_FA/player0

# Inspecter les proprietes
busctl introspect org.bluez \
  /org/bluez/hci0/dev_2C_BE_EB_7F_2D_FA/player0

# Verifier oFono (appels HFP)
busctl call org.ofono / org.ofono.Manager GetModems
# Doit lister un modem avec org.ofono.VoiceCallManager
\end{lstlisting}

%% ===========================================================================
\section{Daemon de liaison RFCOMM : \texttt{nav-bind-daemon.sh}}
%% ===========================================================================

\subsection{Rôle et perimetre}

\texttt{nav-bind-daemon.sh} est le seul demon specifique au projet côte Raspberry Pi.
Il maintient en permanence le binding entre \texttt{/dev/rfcomm1} et le service
\texttt{NavJsonService} de l'application Android NavRelay.

Le numero de canal RFCOMM assigne dynamiquement par Android peut varier à chaque
reconnexion ; il ne peut être connu qu'en consultant le registre SDP du telephone
à la volee, d'où la necessite de ce demon.

\subsection{Variables de configuration}

A partir de la version amelioree du daemon, le MAC du telephone est detecte automatiquement
par \texttt{main.c} et ecrit dans le fichier \texttt{/tmp/connected\_phone\_mac.txt}.

\begin{lstlisting}[style=bashcode,caption={nav-bind-daemon.sh --- detection auto du MAC}]
MAC_FILE="/tmp/connected_phone_mac.txt"
if [ -f "$MAC_FILE" ]; then
  PHONE_MAC=$(cat "$MAC_FILE" 2>/dev/null)
else
  PHONE_MAC="2C:BE:EB:7F:2D:FA"  # Fallback si fichier absent
fi
TARGET_SERVICE_NAME="NavJsonService"
\end{lstlisting}

\emph{Avantage:} Plus besoin de modifier manuellement le MAC. Le daemon se reconfigure
automatiquement lors du reconnexion d'un telephone different.

\textbf{Comment ca marche:}
\begin{enumerate}
  \item \texttt{main.c} detecte le telephone connecte via D-Bus (org.bluez.Device1).
  \item \texttt{main.c} extrait l'adresse MAC et l'ecrit dans \texttt{/tmp/connected\_phone\_mac.txt}.
  \item Le daemon relance la lecture du fichier chaque cycle (5 secondes).
  \item Si le MAC change (nouveau telephone), le daemon se reconfigure automatiquement.
\end{enumerate}

\subsection{Logique de binding detaillee}

Le script tourne dans une boucle infinie avec 5 secondes d'attente entre chaque cycle.

\subsubsection{etape 1 : lecture de l'etat courant de rfcomm1}

\begin{lstlisting}[style=bashcode, showstringspaces=false]
CURRENT_LINE=$(rfcomm | awk -v mac="$PHONE_MAC" \
  '$1=="rfcomm1:" && $2==mac {print $0}')
\end{lstlisting}

La sortie de \texttt{rfcomm} ressemble à :
\begin{verbatim}
rfcomm1: 2C:BE:EB:7F:2D:FA channel 7 config
\end{verbatim}

Si aucune ligne ne correspond, \texttt{CURRENT\_LINE} est vide et un rebind
sera declenche.

\subsubsection{etape 2 : consultation SDP du telephone}

\begin{lstlisting}[style=bashcode, showstringspaces=false]
SDP_OUTPUT=$(sdptool browse "$PHONE_MAC" 2>/dev/null || true)

TARGET_CHANNEL=$(printf '%s\n' "$SDP_OUTPUT" \
  | awk -v name="$TARGET_SERVICE_NAME" '
    $0 ~ "Service Name: "name { in_block=1 }
    in_block && /Channel:/ { print $2; exit }
  ')
\end{lstlisting}

\texttt{sdptool browse} parcourt tous les services Bluetooth du telephone. Si
\texttt{TARGET\_CHANNEL} est vide (NavRelay non demarre ou telephone hors portee),
le cycle recommence apres 5~s.

\subsubsection{etape 3 : decision de rebinding}

Trois cas declenchent un rebind :

\begin{itemize}
  \item \texttt{rfcomm1} n'existe pas ou n'est pas associe à \texttt{PHONE\_MAC}.
  \item \texttt{rfcomm1} existe mais est en etat \texttt{closed}.
  \item \texttt{rfcomm1} existe mais pointe sur un canal different de
    \texttt{TARGET\_CHANNEL}.
\end{itemize}

\subsubsection{etape 4 : rebinding}

\begin{lstlisting}[style=bashcode]
rfcomm release 1 2>/dev/null || true
rfcomm bind 1 "$PHONE_MAC" "$TARGET_CHANNEL"
\end{lstlisting}

Apres un bind reussi, \texttt{/dev/rfcomm1} est disponible et
\texttt{nav\_thread\_fn()} peut l'ouvrir.

\subsection{Installation comme service systemd}

\begin{lstlisting}[style=bashcode]
sudo nano /etc/systemd/system/nav-bind.service
\end{lstlisting}

\begin{lstlisting}[style=confcode,caption={/etc/systemd/system/nav-bind.service}]
[Unit]
Description=RFCOMM binding daemon for NavRelay navigation
After=bluetooth.target
Wants=bluetooth.target

[Service]
Type=simple
ExecStart=/bin/bash /home/pi5/dashboard/nav-bind-daemon.sh
Restart=on-failure
RestartSec=5
StandardOutput=journal
StandardError=journal

[Install]
WantedBy=multi-user.target
\end{lstlisting}

\begin{lstlisting}[style=bashcode]
# 1. Copier le script du daemon
sudo cp nav-bind-daemon.sh /usr/local/bin/nav-bind-daemon.sh

# 2. Rendre executable
sudo chmod +x /usr/local/bin/nav-bind-daemon.sh

# 3. Creer le service systemd
sudo tee /etc/systemd/system/nav-bind.service > /dev/null << 'EOF'
[Unit]
Description=RFCOMM auto-bind daemon for NavJsonService
After=bluetooth.target
Wants=bluetooth.target

[Service]
Type=simple
ExecStart=/usr/local/bin/nav-bind-daemon.sh
Restart=always
RestartSec=5

[Install]
WantedBy=multi-user.target
EOF

# 4. Recharger systemd
sudo systemctl daemon-reload

# 5. Activer au demarrage
sudo systemctl enable nav-bind.service

# 6. Demarrer maintenant
sudo systemctl start nav-bind.service

# 7. Verifier
systemctl status nav-bind.service
\end{lstlisting}

\subsection{Diagnostic et depannage}

\begin{description}
  \item[\texttt{sdptool returned nothing}]
    Le telephone n'est pas connecte. Verifier :
    \texttt{bluetoothctl info 2C:BE:EB:7F:2D:FA} (champ \texttt{Connected: yes}).
  \item[\texttt{cannot find service "NavJsonService"}]
    NavRelay n'est pas demarre. Ouvrir NavRelay et appuyer sur \og Start \fg{}.
  \item[\texttt{rfcomm bind failed}]
    Module noyau manquant : \texttt{sudo modprobe rfcomm}.
    Verifier aussi que le service tourne en root.
  \item[\texttt{/dev/rfcomm1} introuvable apres binding]
    Ajouter l'utilisateur au groupe \texttt{dialout} :
    \texttt{sudo usermod -aG dialout pi5}, puis se reconnecter.
  \item[rfcomm1 se rebind en boucle]
    Qualite de signal Bluetooth insuffisante ou telephone non marque comme
    de confiance (\texttt{trust} dans \texttt{bluetoothctl}).
\end{description}

%% ===========================================================================
\section{Application Android \texttt{NavRelay}}
%% ===========================================================================

\subsection{Structure du projet}

\begin{verbatim}
NavRelay/
  app/src/main/
    AndroidManifest.xml
    java/com/example/navrelay/
      MainActivity.kt
      NavServer.kt
      MapsNotificationListener.kt
    res/layout/
      activity_main.xml
\end{verbatim}

\subsection{Compiler et installer}

\begin{enumerate}
  \item Ouvrir Android Studio et importer le projet \texttt{NavRelay/}.
  \item Verifier que \texttt{compileSdkVersion} $\geq 33$.
  \item Connecter le telephone en USB avec le debogage active.
  \item Cliquer sur \emph{Run} ou : \texttt{./gradlew installDebug}
\end{enumerate}

\subsection{Permissions et manifeste}

\begin{itemize}
  \item \texttt{BLUETOOTH}, \texttt{BLUETOOTH\_ADMIN}, \texttt{BLUETOOTH\_CONNECT},
    \texttt{BLUETOOTH\_ADVERTISE}, \texttt{BLUETOOTH\_SCAN}
  \item \texttt{ACCESS\_FINE\_LOCATION}, \texttt{ACCESS\_COARSE\_LOCATION}
  \item \texttt{POST\_NOTIFICATIONS} (Android~13+)
  \item Service \texttt{MapsNotificationListener} declare avec l'action
    \texttt{BIND\_NOTIFICATION\_LISTENER\_SERVICE}
\end{itemize}

\textbf{IMPORTANT --- etape obligatoire apres installation :} Accorder l'acces aux notifications
manuellement apres le premier lancement :
\begin{enumerate}
  \item Ouvrir \textit{Parametres $\to$ Applications}
  \item Chercher \textit{Acces speciaux} ou \textit{Permissions avancees}
  \item Selectionner \textit{Acces aux notifications}
  \item Activer \texttt{NavRelay}
\end{enumerate}

\texttt{MainActivity} ouvre automatiquement cet ecran via
\texttt{Settings.ACTION\_NOTIFICATION\_LISTENER\_SETTINGS}.

\textbf{Note :} Sur certaines ROM Android, desactiver l'optimisation de batterie
pour NavRelay afin d'eviter que le service soit tue en arriere-plan.

\subsection{\texttt{MainActivity} : demarrage du serveur}

\begin{lstlisting}[style=ktcode,caption={MainActivity.kt --- bouton Start}]
startButton.setOnClickListener {
    NavServer.start(this)
    statusText.text = "NavServer demarre.\n" +
        "Active le service de notifications pour NavRelay."
    startActivity(
        Intent(Settings.ACTION_NOTIFICATION_LISTENER_SETTINGS))
}
\end{lstlisting}

Les permissions Bluetooth et localisation sont demandees au lancement via
\texttt{ActivityCompat.requestPermissions()}, avec les variantes
\texttt{BLUETOOTH\_CONNECT / ADVERTISE / SCAN} pour Android~12+.

\subsection{\texttt{NavServer} : serveur Bluetooth RFCOMM}

\subsubsection{UUID et nom de service}

\begin{lstlisting}[style=ktcode,caption={NavServer.kt --- constantes}]
private const val SERVICE_NAME = "NavJsonService"
val SPP_UUID: UUID =
    UUID.fromString("00001101-0000-1000-8000-00805F9B34FB")
\end{lstlisting}

L'UUID \texttt{00001101-...} est l'UUID standardise du \emph{Serial Port Profile} (SPP).
\texttt{NavJsonService} est le nom que \texttt{nav-bind-daemon.sh} recherche dans le
registre SDP. Ces deux valeurs doivent rester coherentes.

\subsubsection{Cycle de vie du serveur}

\texttt{start()} lance un thread qui :

\begin{enumerate}
  \item Recupere l'adaptateur Bluetooth et verifie qu'il est actif.
  \item Appelle \texttt{listenUsingRfcommWithServiceRecord(SERVICE\_NAME, SPP\_UUID)},
    enregistrant le service dans le registre SDP avec un canal assigne dynamiquement.
    C'est ce canal que le daemon decouvre via \texttt{sdptool}.
  \item Boucle sur \texttt{accept()} pour accepter la connexion du Raspberry Pi.
    À chaque reconnexion, l'ancienne socket et son flux sont fermes et remplaces.
\end{enumerate}

\subsubsection{Envoi des messages JSON}

\begin{lstlisting}[style=ktcode,caption={NavServer.kt --- send()}]
fun send(json: JSONObject) {
    val out = outRef.get() ?: return
    try {
        val line = json.toString() + "\n"
        out.write(line.toByteArray(Charsets.UTF_8))
        out.flush()
    } catch (e: Exception) {
        e.printStackTrace()
        stop()
    }
}
\end{lstlisting}

Chaque message est une ligne JSON terminee par \verb|\n| en UTF-8. En cas d'exception
d'ecriture (deconnexion Bluetooth), le serveur s'arrête proprement et attend une
nouvelle connexion.

\subsubsection{Format JSON}

\begin{lstlisting}[style=confcode,caption={Exemple de message JSON envoye par NavRelay}]
{
  "timestamp": 1711550400,
  "next_action": "Tournez a droite sur Rue de la Paix",
  "distance_to_next_m": 320,
  "speed_kmh": 48.5,
  "lat": 48.8566,
  "lon": 2.3522
}
\end{lstlisting}

\subsection{\texttt{MapsNotificationListener}}

\subsubsection{Interception des notifications}

\texttt{onNotificationPosted()} filtre les notifications dont le \texttt{packageName}
est \texttt{"com.google.android.apps.maps"} et extrait :

\begin{itemize}
  \item \texttt{EXTRA\_TITLE} : distance jusqu'à la prochaine manœuvre (ex : \og 320~m \fg{}).
  \item \texttt{EXTRA\_TEXT} : instruction textuelle (ex : \og Tournez à droite \fg{}).
  \item \texttt{EXTRA\_SUB\_TEXT} : information complementaire.
\end{itemize}

Le champ \texttt{next\_action} est determine par ordre de priorite :
\texttt{EXTRA\_TEXT} $>$ \texttt{EXTRA\_SUB\_TEXT} $>$ \texttt{EXTRA\_TITLE}.

\subsubsection{Extraction de la distance}

\texttt{parseDistanceMeters()} analyse la chaîne concatenee avec une expression
reguliere reconnaissant les formats \og \texttt{320 m} \fg{} et \og \texttt{1.2 km} \fg{}.
Les kilometres sont convertis en metres pour normaliser \texttt{distance\_to\_next\_m}.

\subsubsection{Donnees GPS}

\texttt{requestLocationUpdates()} demande des mises à jour GPS toutes les secondes
(deplacement minimal 1~m). La vitesse \texttt{Location.speed} (m/s) est convertie
en km/h (\texttt{speed * 3.6f}) dans \texttt{onLocationChanged()}, puis integree
dans le prochain message JSON.

\subsubsection{Verifier la connexion depuis le Pi}

\begin{lstlisting}[style=bashcode]
# Verifier que NavJsonService est visible en SDP
sdptool browse 2C:BE:EB:7F:2D:FA | grep -A5 "NavJsonService"
# Sortie attendue :
#   Service Name: NavJsonService
#   Channel: 7   (numero variable)

# Lire le flux JSON en direct
cat /dev/rfcomm1
\end{lstlisting}

%% ===========================================================================
\section{Application LVGL (\texttt{lv\_port\_vscode/src})}
%% ===========================================================================

\subsection{Macros de configuration}

\begin{lstlisting}[style=ccode,caption={main.c --- macros de configuration}]
#define SCREEN_W              1080
#define SCREEN_H              1080
#define MUSIC_COVER_PATH_FS   "/home/pi5/dashboard/cover.jpg"
#define MUSIC_COVER_PATH      "A:/home/pi5/dashboard/cover.jpg"
#define CENTER_RADIUS         250
#define BOTTOM_PLAYER_HEIGHT  250
\end{lstlisting}

\texttt{MUSIC\_COVER\_PATH\_FS} est le chemin POSIX reel pour ecrire et verifier la
jaquette. \texttt{MUSIC\_COVER\_PATH} utilise le prefixe \texttt{A:} du pilote LVGL POSIX
initialise par \texttt{lv\_fs\_posix\_init()}.

\subsection{Structures de donnees partagees}

Toutes les donnees partagees entre les threads et le timer UI sont stockees dans
l'instance globale \texttt{g\_state} de type \texttt{app\_state\_t}, protegee par
\texttt{g\_state.lock} (mutex POSIX).

\begin{lstlisting}[style=ccode,caption={main.c --- structures de donnees}]
typedef struct {
    char   instruction[256]; /* Instruction de virage courante */
    double distance;         /* Distance en metres             */
    double next_distance;
    char   street_name[256]; /* Nom de la rue courante         */
    double speed;            /* Vitesse vehicule en km/h       */
    double lat;
    double lon;
} nav_info_t;

typedef struct {
    char         title[256];
    char         artist[256];
    char         album[256];
    char         status[32];     /* "playing", "paused", etc.  */
    unsigned int position_ms;
    unsigned int duration_ms;
    char         cover_path[512];
} music_info_t;

typedef struct {
    char state[32];       /* "active","incoming","dialing"...  */
    char line_id[64];     /* Numero de telephone               */
    char name[64];        /* Nom du contact                    */
    char call_path[256];  /* Chemin D-Bus de l'appel oFono     */
} call_info_t;

typedef struct {
    nav_info_t      nav;
    music_info_t    music;
    call_info_t     current_call;
    int             has_incoming_call;
    pthread_mutex_t lock;
} app_state_t;
\end{lstlisting}

\subsection{Thread de navigation : \texttt{nav\_thread\_fn()}}

\subsubsection{Connexion et reconnexion à \texttt{/dev/rfcomm1}}

\begin{lstlisting}[style=ccode,caption={main.c --- boucle de reconnexion rfcomm1}]
while (1) {
    if (fd < 0) {
        fd = open("/dev/rfcomm1", O_RDONLY | O_NOCTTY);
        if (fd < 0) { sleep(2); continue; }
        f = fdopen(fd, "r");
        if (!f) { close(fd); fd = -1; sleep(2); continue; }
    }
    if (fgets(line, sizeof(line), f)) {
        /* strip \r\n, parse JSON, mise a jour g_state.nav */
    } else {
        /* EOF ou erreur : fermeture et reconnexion */
        fclose(f); f = NULL; fd = -1; sleep(2);
    }
}
\end{lstlisting}

Le port est ouvert avec \texttt{O\_NOCTTY} pour eviter d'en faire le terminal de
contrôle du processus.

\subsubsection{Parsing JSON}

\texttt{parse\_nav\_json()} utilise \texttt{json\_tokener\_parse()} de \texttt{json-c}.
Elle extrait \texttt{next\_action}, \texttt{distance\_to\_next\_m}, \texttt{street\_name},
\texttt{speed\_kmh}, \texttt{lat}, \texttt{lon} et remplit \texttt{nav\_info\_t}.
Si \texttt{street\_name} est absent, l'UI affiche \og Unknown Street \fg{}.

\subsection{Thread musique : \texttt{music\_thread\_fn()}}

\subsubsection{Connexion D-Bus et decouverte du lecteur}

\texttt{dbus\_connect\_system()} initialise une connexion au bus D-Bus systeme via
\texttt{dbus\_bus\_get(DBUS\_BUS\_SYSTEM, \&err)}.

\texttt{find\_player\_path()} envoie \texttt{GetManagedObjects} à \texttt{org.bluez}
sur \texttt{/} et cherche le premier objet exposant \texttt{org.bluez.MediaPlayer1}.
Le chemin resultant est de la forme :

\begin{verbatim}
/org/bluez/hci0/dev_2C_BE_EB_7F_2D_FA/player0
\end{verbatim}

\subsubsection{Lecture des proprietes}

Toutes les 2 secondes, \texttt{get\_music\_properties()} appelle
\texttt{GetAll("org.bluez.MediaPlayer1")} et extrait :

\begin{itemize}
  \item \texttt{Status} (chaîne) $\to$ \texttt{g\_state.music.status}
  \item \texttt{Position} (uint32, ms) $\to$ \texttt{g\_state.music.position\_ms}
  \item \texttt{Track} (dict) $\to$ \texttt{Title}, \texttt{Artist}, \texttt{Album},
    \texttt{Duration}
\end{itemize}

En cas de changement de piste (couple titre+artiste), \texttt{download\_itunes\_cover()}
est appelee pour recuperer la nouvelle jaquette.

\subsubsection{Commandes de contrôle}

\texttt{send\_player\_method(method)} envoie \texttt{Play}, \texttt{Pause},
\texttt{Next} ou \texttt{Previous} à \texttt{org.bluez.MediaPlayer1} avec un
timeout de 5000~ms. Ces appels sont declenches par les boutons LVGL sous le
verrou \texttt{g\_state.lock}.

\subsection{Thread appels : \texttt{calls\_thread\_fn()}}

\subsubsection{Decouverte du modem HFP}

\texttt{find\_modem\_with\_voicecall()} interroge
\texttt{org.ofono.Manager.GetModems()} et selectionne le modem dont la liste
d'interfaces contient \texttt{"org.ofono.VoiceCallManager"} :

\begin{verbatim}
/hfp/org/bluez/hci0/dev_2C_BE_EB_7F_2D_FA
\end{verbatim}

\subsubsection{Polling des appels}

Toutes les 2 secondes, \texttt{get\_calls()} envoie
\texttt{VoiceCallManager.GetCalls()} et extrait pour chaque appel \texttt{State},
\texttt{LineIdentification} et \texttt{Name}. Le thread met à jour
\texttt{has\_incoming\_call} si un appel est en etat \texttt{"incoming"} ou
\texttt{"waiting"}.

\subsubsection{Decrocher et raccrocher}

\texttt{call\_voicecall\_method(call\_path, method)} envoie \texttt{Answer} ou
\texttt{Hangup} à \texttt{org.ofono.VoiceCall} (timeout 5~s).

\subsection{Timer UI (100~ms)}

\begin{lstlisting}[style=ccode,caption={main.c --- callback du timer LVGL}]
static void ui_update_timer(lv_timer_t *timer) {
    update_nav_display();
    update_music_display();
    update_call_display();
    speedo_set_speed((int32_t)g_state.nav.speed);
}
\end{lstlisting}

\begin{description}
  \item[\texttt{update\_nav\_display()}]
    Instruction, distance (m ou km), nom de rue ou \og Unknown Street \fg{}.
  \item[\texttt{update\_music\_display()}]
    Titre, artiste, jaquette rechargee via \texttt{stat()} +
    \texttt{lv\_image\_set\_src()}, label Play/Pause.
  \item[\texttt{update\_call\_display()}]
    Affiche/masque le panneau (\texttt{LV\_OBJ\_FLAG\_HIDDEN}), nom ou
    \og Unknown \fg{} et numero de telephone.
  \item[\texttt{speedo\_set\_speed()}]
    Compteur mis à jour avec \texttt{g\_state.nav.speed} en km/h entier.
\end{description}

\subsection{Layout LVGL}

\begin{center}
\begin{tabular}{llll}
\toprule
\textbf{Widget} & \textbf{Taille} & \textbf{Position} & \textbf{Fond} \\
\midrule
Compteur (speedo) & $1080\times1080$ & Plein ecran & transparent \\
Conteneur nav. & $500\times500$ & Centre, +100~px Y & \texttt{\#1a1a2e} \\
Conteneur musique & $830\times220$ & Bas, $-150$~px Y & \texttt{\#0f0f1e} \\
Jaquette album & $100\times100$ & Gauche musique & --- \\
Panneau appel & $600\times250$ & Par-dessus musique & \texttt{\#393a4f} \\
Pair Button & Petit & Sous musique & Bleu \\
\bottomrule
\end{tabular}
\end{center}

Le panneau d'appel est cree avec \texttt{LV\_OBJ\_FLAG\_HIDDEN} et n'est rendu
visible que lors d'un appel entrant. Le Pair Button appelle \texttt{bluetooth\_pair\_device\_start()}.

\subsection{Callbacks des boutons}

Six boutons LVGL sont definis :

\begin{center}
\begin{tabular}{lll}
\toprule
\textbf{Bouton} & \textbf{Fonction} & \textbf{Interface} \\
\midrule
Pair & \texttt{bluetooth\_pair\_device\_start()} & \texttt{bt\_make\_discoverable.sh} \\
Precedent & \texttt{Previous} & \texttt{org.bluez.MediaPlayer1} \\
Play/Pause & \texttt{Play} / \texttt{Pause} & \texttt{org.bluez.MediaPlayer1} \\
Suivant & \texttt{Next} & \texttt{org.bluez.MediaPlayer1} \\
Decrocher & \texttt{Answer} & \texttt{org.ofono.VoiceCall} \\
Raccrocher & \texttt{Hangup} & \texttt{org.ofono.VoiceCall} \\
\bottomrule
\end{tabular}
\end{center}

%% ===========================================================================
\section{Compteur de vitesse LVGL : \texttt{speedo\_interactive.c/.h}}
%% ===========================================================================

\subsection{API publique}

\begin{lstlisting}[style=ccode,caption={speedo\_interactive.h}]
void speedo_interactive_create(void);
void speedo_set_speed(int32_t speed);
\end{lstlisting}

\texttt{speedo\_interactive\_create()} doit être appelee une seule fois dans
\texttt{create\_ui()}. \texttt{speedo\_set\_speed()} est appelee depuis le timer
UI à chaque tick de 100~ms.

\subsection{Parametres du cadran}

\begin{lstlisting}[style=ccode,caption={speedo\_interactive.c --- constantes}]
#define SPEEDO_W    1080    /* Largeur totale (px)               */
#define SPEEDO_H    1080    /* Hauteur totale (px)               */
#define SPEED_MIN   0       /* Vitesse minimale (km/h)           */
#define SPEED_MAX   260     /* Vitesse maximale (km/h)           */
#define ARC_SWEEP   180     /* Balayage angulaire total (degres) */
#define ARC_START   270     /* Angle de depart (bas du cercle)   */
\end{lstlisting}

\subsection{Conversion vitesse/angle et aiguille}

\texttt{speed\_to\_angle()} calcule l'angle par interpolation lineaire :

\[
  \theta = \texttt{ARC\_START} + \frac{v - v_{\min}}{v_{\max} - v_{\min}}
    \times \texttt{ARC\_SWEEP}
\]

\texttt{update\_needle()} convertit $\theta$ en radians, calcule les deux extremites
de l'aiguille sur la couronne externe du cadran et met à jour l'objet ligne LVGL
\texttt{needle\_line}.

\texttt{create\_ticks\_and\_labels()} cree pour chaque pas de 5~km/h :
\begin{itemize}
  \item Graduation majeure tous les 10~km/h (plus longue, couleur claire).
  \item Graduation mineure tous les 5~km/h.
  \item Label numerique (police \texttt{lv\_font\_montserrat\_18}) tous les 25~km/h.
\end{itemize}

\texttt{speedo\_interactive\_create()} cree les graduations, l'aiguille, et deux
labels centraux : la valeur numerique (police 48~pt) et l'unite \og km/h \fg{}.

%% ===========================================================================
\section{Recuperation de jaquette : \texttt{album\_art.c}}
%% ===========================================================================

\subsection{Signature}

\begin{lstlisting}[style=ccode]
int download_itunes_cover(const char *artist,
                          const char *album,
                          const char *out_path);
\end{lstlisting}

Dependances : \texttt{libcurl} et \texttt{json-c}. Retourne 1 en succes, 0 sinon.

\subsection{Flux d'execution}

\begin{enumerate}
  \item Encodage URL du terme (\textit{artiste} + espace + \textit{album}) via
    \texttt{curl\_easy\_escape()}.
  \item Requête GET vers :
    \url{https://itunes.apple.com/search?term=<query>&entity=album&limit=1}
  \item Accumulation de la reponse JSON en memoire via \texttt{write\_to\_memory()}
    (structure \texttt{Memory} avec \texttt{realloc()} dynamique).
  \item Extraction de \texttt{"artworkUrl100"} dans \texttt{results[0]}.
  \item Telechargement de l'image vers \texttt{out\_path} via
    \texttt{write\_to\_file()}.
\end{enumerate}

Si le telechargement echoue (pas de reseau, artiste inconnu), la fonction retourne 0
et l'ancienne jaquette reste affichee jusqu'au prochain changement de piste.

%% ===========================================================================
\section{Ordre de demarrage complet}
%% ===========================================================================

\begin{enumerate}
  \item \textbf{Demarrer le Raspberry Pi} : \texttt{bluetooth}, \texttt{ofono} et
    \texttt{nav-bind} demarrent automatiquement via systemd (apres SETUP.sh).
  \item \textbf{Sur le telephone} : activer le Bluetooth.
  \item \textbf{Lancer le dashboard} : \texttt{./bin/main} depuis \texttt{lv\_port\_vscode/build}.
  \item \textbf{Cliquer sur le bouton Pair} dans le dashboard.
  \item \textbf{Appairer le telephone} via Parametres Bluetooth du Pi.
  \item \textbf{Ouvrir NavRelay} et appuyer sur \textbf{Start}.
  \item \textbf{Lancer Google Maps} et une navigation.
\end{enumerate}

\begin{center}
\begin{tabular}{lll}
\toprule
\textbf{Thread} & \textbf{Source de donnees} & \textbf{Intervalle} \\
\midrule
\texttt{nav\_thread\_fn} & \texttt{/dev/rfcomm1} (JSON) & Bloquant sur \texttt{fgets()} \\
\texttt{music\_thread\_fn} & D-Bus \texttt{MediaPlayer1} & 2 secondes \\
\texttt{calls\_thread\_fn} & D-Bus \texttt{VoiceCallManager} & 2 secondes \\
Timer UI & \texttt{lv\_timer\_handler()} & 100 ms \\
\bottomrule
\end{tabular}
\end{center}

%% ===========================================================================
\section{Verification et validation pre-lancement}
%% ===========================================================================

\subsection{Checklist systeme (apres SETUP.sh)}

\begin{itemize}
  \item[\checkmark] Bluetooth demarre : \texttt{sudo systemctl status bluetooth}
  \item[\checkmark] oFono demarre : \texttt{sudo systemctl status ofono}  
  \item[\checkmark] nav-bind service actif : \texttt{sudo systemctl status nav-bind.service}
  \item[\checkmark] Adaptateur visible : \texttt{hciconfig} (affiche \texttt{hci0})
\end{itemize}

\subsection{Checklist appairage}

\begin{itemize}
  \item[\checkmark] Telephone detectable : \texttt{bluetoothctl scan on}
  \item[\checkmark] Telephone appaire : \texttt{bluetoothctl paired-devices}
  \item[\checkmark] Telephone de confiance : \texttt{bluetoothctl trust <MAC>}
  \item[\checkmark] Telephone connecte : \texttt{bluetoothctl info <MAC>} → \texttt{Connected: yes}
\end{itemize}

\subsection{Checklist application}

\begin{itemize}
  \item[\checkmark] Executable compile : \texttt{ls -la lv\_port\_vscode/bin/main}
  \item[\checkmark] Affichage X11 disponible : \texttt{echo \$DISPLAY} → \texttt{:0}
  \item[\checkmark] Permissions dialout : \texttt{groups pi5} → contient \texttt{dialout}
\end{itemize}

%% ===========================================================================
\section{Depannage general}
%% ===========================================================================

\subsection{\og Waiting for route...\fg{} en permanence}

\begin{itemize}
  \item Verifier que \texttt{/dev/rfcomm1} existe : \texttt{ls /dev/rfcomm*}
  \item Tester la reception : \texttt{cat /dev/rfcomm1} (doit afficher des lignes JSON).
  \item Verifier les permissions : \texttt{sudo usermod -aG dialout pi5}.
  \item Dans NavRelay : cliquer sur \textbf{Start}, verifier que l'etat passe à \texttt{"NavServer running"}.
\end{itemize}

\subsection{\og No MediaPlayer\fg{}}

\begin{itemize}
  \item Verifier A2DP avec \texttt{busctl tree org.bluez} (chercher \texttt{/player0}).
  \item Jouer un morceau sur le telephone pour que BlueZ cree \texttt{MediaPlayer1}.
  \item Redemarrer le dashboard si le player est apparu apres son demarrage.
\end{itemize}

\subsection{Aucun appel detecte}

\begin{itemize}
  \item Verifier oFono : \texttt{systemctl status ofono}.
  \item Verifier \texttt{Disable=Headset} dans \texttt{/etc/bluetooth/main.conf}.
  \item Verifier l'existence de \texttt{/etc/pipewire/pipewire-pulse.conf.d/10-ofono.conf}.
  \item Tester : \texttt{busctl call org.ofono / org.ofono.Manager GetModems}.
\end{itemize}

\subsection{La jaquette ne s'affiche pas}

\begin{itemize}
  \item Tester la connectivite :
    \texttt{curl "https://itunes.apple.com/search?term=test\&entity=album\&limit=1"}
  \item Verifier que \texttt{/home/pi5/dashboard/} est accessible en ecriture.
  \item S'assurer que \texttt{libcurl} et \texttt{libjson-c} sont lies à la compilation.
\end{itemize}

\subsection{Pair Button ne fonctionne pas}

\begin{itemize}
  \item Verifier que \texttt{bt\_make\_discoverable.sh} existe et est executable :
    \texttt{ls -la /home/pi5/dashboard/bt\_make\_discoverable.sh}
  \item Tester manuellement : \texttt{bash /home/pi5/dashboard/bt\_make\_discoverable.sh}
  \item Verifier les logs dashboard : \texttt{journalctl -fu dashboard.service}
\end{itemize}

\subsection{Commandes de logs utiles}

\begin{lstlisting}[style=bashcode]
journalctl -fu nav-bind.service        # Daemon RFCOMM
journalctl -fu dashboard.service       # Application LVGL
journalctl -fu bluetooth.service       # BlueZ
journalctl -fu ofono.service           # oFono
dbus-monitor --system \
  "type=signal,sender=org.bluez"       # Evenements BlueZ
dbus-monitor --system \
  "type=signal,sender=org.ofono"       # Evenements oFono
\end{lstlisting}

%% ===========================================================================
\section{Limitations connues}
%% ===========================================================================

\begin{description}
  \item[Nom de l'appelant (HFP)]
    La propriete \texttt{Name} de \texttt{org.ofono.VoiceCall} est generalement vide
    avec les telephones Android modernes. Seul le numero (\texttt{LineIdentification})
    est affiche de maniere fiable.
  \item[Jaquette Bluetooth (AVRCP Cover Art)]
    BlueZ ne supporte pas AVRCP~1.6 de maniere stable. La jaquette est recuperee via
    l'API iTunes Search, ce qui necessite une connexion Internet et introduit un
    delai au changement de piste.
  \item[Redecouverte du player BlueZ]
    \texttt{find\_player\_path()} est appelee une seule fois au demarrage du thread
    musique. En cas de deconnexion Bluetooth pendant l'execution, il faut redemarrer
    le dashboard pour redecouvrir le player.
  \item[Canal RFCOMM variable]
    Le canal assigne à \texttt{NavJsonService} par Android peut changer à chaque
    redemarrage. Ce cas est gere automatiquement par \texttt{nav-bind-daemon.sh}
    via le registre SDP.
  \item[Approach systeme vs. script]
    La nouvelle approche du Pair Button utilise l'interface Bluetooth native du Pi OS
    plutôt quune tentative automatisee, ce qui elimine les conflits logiciels.
\end{description}

%% ===========================================================================
\section{Fichiers du projet}
%% ===========================================================================

\subsection{Fichiers essentiels}

\begin{description}
  \item[\texttt{SETUP.sh}]
    Script d'installation une fois (4 KB). À executer apres la premiere installation,
    installe tous les paquets et services systemd.
  \item[\texttt{QUICK\_START.md}]
    Guide d'utilisation pour l'utilisateur final. À consulter pour le demarrage
    quotidien et les procedures courantes.
  \item[\texttt{nav-bind-daemon.sh}]
    Daemon de liaison RFCOMM. Maintient le binding entre \texttt{/dev/rfcomm1} et
    le service NavRelay du telephone. À adapter pour un autre MAC.
  \item[\texttt{bt\_make\_discoverable.sh}]
    Rend le Pi decouvrable pendant 3 minutes. Appele par le Pair Button.
  \item[\texttt{bt\_setup\_audio.sh}]
    Configuration optionnelle de PulseAudio pour Bluetooth. Generalement non necessaire
    apres SETUP.sh.
  \item[\texttt{cover.jpg}]
    Cache pour la jaquette de l'album courant (vide initialement).
\end{description}

\end{document}
